\documentclass[11pt]{article}

\usepackage[margin=1in]{geometry}
\usepackage{amsmath,amssymb}
\usepackage{graphicx}
\usepackage{booktabs}
\usepackage{hyperref}
\usepackage{cite}
\usepackage{listings}
\usepackage{xcolor}

\hypersetup{
    colorlinks=true,
    linkcolor=blue,
    citecolor=blue,
    urlcolor=blue,
}

\lstset{
    basicstyle=\ttfamily\small,
    breaklines=true,
    frame=single,
    language=Python,
    numbers=left,
    xleftmargin=1em,
}

\title{\bf Dense Reward Shaping Eliminates Policy Memorization\\
in PPO Agents on Atari Breakout%
\thanks{Code and reproduction: \url{https://github.com/mharrell/breakout-reactive-ppo}.
Full project history: \url{https://github.com/mharrell/BreakoutBot}.
Split-watcher demonstration: \url{https://www.youtube.com/watch?v=6ixVwQm7u5Y}.}}

\author{Mike Harrell\\
\texttt{mikey.harrell@gmail.com}}

\date{August 2026}

\begin{document}

\maketitle

\begin{abstract}
Proximal Policy Optimization (PPO) agents trained on deterministic Atari environments
routinely converge to memorized action sequences---fixed scripts that replay the same
buttons in the same order regardless of game state---rather than developing reactive,
ball-tracking policies. Despite the widespread use of sticky actions as a standard
mitigation, we present evidence from 124 controlled experiments showing that no
environment-modification approach---sticky actions, entropy tuning, cursor
wrappers, dynamics randomization, adversarial obstacles, or frame
skipping---produces a genuinely reactive argmax on Atari Breakout. Every
diagnostic in the standard RL toolbox measures the policy distribution, not the
argmax; the models appear reactive while evaluating to deterministic scripts.
We introduce the \emph{split-watcher}, a diagnostic that measures the argmax
directly by running the same model on two different brick layouts with
independent per-side predictions, and show that it detects memorization where
distribution-level diagnostics fail. We then demonstrate that a simple dense
proximity reward---three lines of code that reward the paddle for horizontal
closeness to the ball during descent---produces the first verified reactive PPO
argmax on Atari Breakout. The model achieves 0/240 perfect transfers on the
split-watcher (100\% no-timing ALT retention), an intervention gradient AUC of 0.421,
and a stochastic best score of 216 on clean evaluation. The key insight is that
reward shaping changes what the optimizer considers optimal, rather than trying
to make memorized behavior non-viable through environment engineering.
\end{abstract}

%-------------------------------------------------------------------------------
\section{Introduction}

Proximal Policy Optimization (PPO)~\cite{schulman2017ppo} is one of the most
widely used deep reinforcement learning algorithms, and the Arcade Learning
Environment (ALE)~\cite{bellemare2013ale} remains a standard benchmark. However,
a well-known issue with training on deterministic ALE games is \emph{policy
memorization}: rather than learning to react to visual input, the agent learns
a fixed action sequence---a ``script''---that exploits deterministic environment
dynamics to achieve a consistent score without any genuine perception of game
state~\cite{machado2018revisiting,zhang2018overfitting}.

The standard mitigation is sticky actions (repeat\_action\_probability=0.25),
proposed by Machado et al.~\cite{machado2018revisiting} and widely adopted as
part of the ALE evaluation protocol. However, Zhang et al.~\cite{zhang2018overfitting}
showed that sticky actions do not prevent memorization in deep ConvNet-based RL
agents: convolutional networks are ``automatically robust'' to 25\% action noise,
and a memorized sequence survives with only minor score degradation.

This paper reports the results of a systematic investigation into policy
memorization in PPO agents on Atari Breakout. Over 124 controlled experiments
spanning every major proposed mitigation, we found that \emph{no}
environment-modification approach produced a reactive argmax. Every model
that appeared reactive under standard diagnostics was revealed as a memorized
script when the argmax---rather than the policy distribution---was tested
directly.

The contribution of this work is threefold:
\begin{enumerate}
    \item We identify a fundamental confound in the standard RL diagnostic
    toolkit: every commonly used metric measures the \emph{policy distribution},
    not the \emph{argmax action}, and distribution-level reactivity routinely
    coexists with argmax memorization.
    \item We introduce the \emph{split-watcher}, a simple diagnostic that
    measures argmax behavior directly by running the same model on two
    different initial states with independent per-side predictions.
    \item We show that a three-line dense proximity reward---directly rewarding
    the paddle for being horizontally close to the descending ball---produces
    the first verified reactive PPO argmax on Atari Breakout, achieving
    0/240 perfect transfers on the split-watcher.
\end{enumerate}

%-------------------------------------------------------------------------------
\section{Related Work}

\subsection{Memorization in Deterministic Environments}

The problem of policy memorization in deterministic RL environments has been
recognized since the early days of deep RL on the ALE. Machado et
al.~\cite{machado2018revisiting} proposed sticky actions
(repeat\_action\_probability=0.25) as a standard mitigation and documented the
prevalence of deterministic exploitation across Atari games. However, Zhang et
al.~\cite{zhang2018overfitting} directly tested whether sticky actions prevent
memorization in DQN and other deep RL agents, finding that ``stochasticity could
neither prevent deep RL agents from serious overfitting nor detect overfitted
agents effectively.'' The present work independently confirms Zhang et al.'s
finding across PPO agents: every sticky-trained model collapsed to a
deterministic script when tested without sticky actions.

Subsequent work has explored various approaches to forcing reactive behavior.
Justesen et al.~\cite{justesen2018illuminating} showed that agents can exploit
deterministic patterns even in procedurally generated environments. Cobbe et
al.~\cite{cobbe2020procgen} demonstrated that Procgen's procedural generation
improves generalization but does not eliminate it. The Go-Explore
algorithm~\cite{ecoffet2021goexplore} achieved superhuman performance on Montezuma's
Revenge but did so through explicit state-space exploration, not reactive
policy learning.

\subsection{Reward Shaping}

Reward shaping---adding auxiliary rewards to guide learning---has a long history
in RL~\cite{ng1999shaping}. In deep RL, auxiliary tasks and dense rewards have
been used to improve sample efficiency and final performance~\cite{jaderberg2017auxiliary}.
The present work differs in motivation: rather than using reward shaping to
accelerate learning, we use it to \emph{change the optimum itself}---making
ball-tracking the reward-maximizing behavior rather than trying to make scripts
non-viable.

\subsection{Diagnosing Memorization}

Standard evaluation protocols for Atari RL rely on aggregate metrics (mean
episode reward, unique score count) that do not distinguish reactive behavior
from memorized scripts. The memorization check used in this project's earlier
work~\footnote{Full project history and methodology at
\url{https://github.com/mharrell/BreakoutBot}} counts unique deterministic
scores across evaluation episodes---but we found that score variance from
life-loss timing and ball-bounce stochasticity produces false MULTIPLE\_SCRIPTS
verdicts for confirmed-memorized models. The split-watcher introduced here is,
to our knowledge, the first diagnostic that directly measures argmax behavior
rather than properties of the policy distribution or score distribution.

%-------------------------------------------------------------------------------
\section{Method}

\subsection{The Split-Watcher Diagnostic}

The split-watcher runs the same trained model on two different Breakout layouts
simultaneously, making independent predictions for each side. The left side
presents a full brick wall; the right side presents an altered layout---half
the bricks cleared, or a different arrangement. Both sides use the same model
weights and the same argmax action selection.

The key insight is that different brick layouts cause different ball
bounces. On the altered side, the ball hits different bricks at different
angles and arrives at the paddle from different positions. A reactive policy
that tracks the ball \emph{must} move differently on the two sides. A memorized
script produces identical action sequences regardless.

We compute the Pearson correlation of paddle X-positions between the two sides:

\begin{equation}
    \text{px\_corr} = \frac{\text{cov}(\mathbf{p}_{\text{full}},
    \mathbf{p}_{\text{alt}})}
    {\sigma_{\text{full}} \sigma_{\text{alt}}}
\end{equation}

A px\_corr $>$ 0.99 with ALT score $\approx$ FULL score is classified as a
\emph{perfect transfer}---definitive memorization, physically impossible for a
reactive policy.

We evaluate across three layout types:
\begin{itemize}
    \item \textbf{RIGHT\_HALF}: bricks 0--17 cleared (right half of the wall)
    \item \textbf{LEFT\_HALF}: bricks 18--35 cleared (left half of the wall)
    \item \textbf{RANDOM\_50}: 50\% of bricks cleared at random
\end{itemize}

We also provide a \emph{no-timing variant} that removes the NoopResetEnv
wrapper (which adds 0--30 random NOOP frames at reset), eliminating timing
offsets as a confound. Both sides start at frame 0 identically; any divergence
comes purely from different brick layouts.

Full code for the split-watcher is available in the companion repository at
\url{https://github.com/mharrell/breakout-reactive-ppo}.

\subsection{Intervention Gradient}

As a complementary diagnostic, we developed the \emph{intervention gradient}
test. Mid-game, after the ball has been in play for sufficient time for the
policy to establish a trajectory, the ball's horizontal position (ALE RAM
address 99) is teleported by a displacement $\Delta x \in \{0, \pm8, \pm15,
\pm30, \pm45, \pm60\}$ pixels. We record whether the paddle's next action
reverses direction to follow the teleported ball. A memorized script ignores
the teleport and continues its fixed sequence. A reactive tracker reverses
direction to track the new ball position.

A center-hold dead script serves as the calibration baseline. The reversal
rate is plotted against displacement magnitude to produce a dose-response
curve; the area under this curve (AUC) quantifies overall reactivity.

\subsection{Proximity Reward Wrapper}

The proximity reward wrapper reads three ALE RAM addresses at each timestep:
paddle X (72), ball X (99), and ball Y (101). When the ball is descending
toward the paddle ($\text{ball\_y} > 100$), a bonus is computed:

\begin{equation}
    \text{bonus} = 0.05 \cdot \max\left(0, \; 1 - \frac{|\text{paddle\_x} -
    \text{ball\_x}|}{80}\right)
\end{equation}

The bonus ranges from 0.0 (paddle and ball separated by $\geq$80 pixels, half
the playfield width) to 0.05 (paddle centered directly under the ball). At
$\sim$2,000 descent frames per game, perfect tracking earns approximately 50
bonus points---roughly 5--10\% of a typical game score and equivalent to 7
yellow bricks. The scale is intentionally small: large enough to shift the
optimum over thousands of frames, small enough not to overwhelm the game's
natural reward signal (1.0--7.0 per brick).

Critically, the proximity bonus is applied \emph{during training only}.
Evaluation and memorization checks use clean Breakout (no wrapper), testing
whether the tracking behavior transfers to the original reward function.

\subsection{Training Configuration}

\begin{table}[h]
\centering
\begin{tabular}{ll}
\toprule
\textbf{Parameter} & \textbf{Value} \\
\midrule
Algorithm & PPO (Stable-Baselines3) \\
Architecture & NatureCNN \\
Environments & 32 (DummyVecEnv) \\
n\_steps & 128 \\
Batch size & 1024 \\
n\_epochs & 4 \\
$\gamma$ & 0.99 \\
Entropy coefficient & 0.006 \\
Learning rate & $2.5\times10^{-4} \rightarrow 10^{-5}$ (linear) \\
Clip range & $0.2 \rightarrow 0.05$ (linear) \\
Total steps & 25,000,000 \\
Environment & ALE/Breakout-v5, frameskip=4, no sticky actions \\
Training wrapper & EpisodicLifeEnv + ProximityRewardWrapper \\
Eval wrapper & EpisodicLifeEnv (clean Breakout, no proximity reward) \\
Hardware & Single NVIDIA RTX 3060 Ti, Windows 11 \\
Training time & $\sim$6.5 hours for 25M steps \\
\bottomrule
\end{tabular}
\caption{Training configuration.}
\label{tab:config}
\end{table}

%-------------------------------------------------------------------------------
\section{Results}

\subsection{Split-Watcher Verification}

\begin{table}[h]
\centering
\begin{tabular}{lcccc}
\toprule
\textbf{Checkpoint} & \textbf{Games} & \textbf{Perfect Transfers} &
\textbf{ALT Retention} & \textbf{Divergence} \\
\midrule
\multicolumn{5}{c}{\textit{No-timing variant (no NoopResetEnv)}} \\
\midrule
best (19.2M) & 60 & 0 & 100\% & 62.4\% \\
final (25M)  & 60 & 0 & 100\% & 62.9\% \\
\midrule
\multicolumn{5}{c}{\textit{With NoopResetEnv (0--30 frame offset)}} \\
\midrule
best (19.2M) & 60 & 0 & 46\% & 71.4\% \\
final (25M)  & 60 & 0 & 59\% & 70.6\% \\
\bottomrule
\end{tabular}
\caption{Split-watcher results for PPO\_124. 0/240 total perfect transfers.}
\label{tab:split}
\end{table}

Table~\ref{tab:split} shows the split-watcher results. Across 240 games (120
no-timing, 120 with NoopResetEnv), the model produced \textbf{zero perfect
transfers}. In the no-timing variant, ALT score retention was 100\%---the model
cleared every brick on every layout, every game. With NoopResetEnv, retention
dropped to 46--59\% due to frame-stack corruption on the first serve
(0--30 random NOOP frames can cause the model to miss the initial ball
trajectory), but the improved retention from 46\% to 59\% between checkpoints
indicates increasing robustness to timing variation.

For comparison, every prior model (PPO\_111--118, cursor generation; BeamRider
baseline and MULTILIFE) showed at least one perfect transfer under the
split-watcher. PPO\_124 is the first model in the project's history to achieve
zero perfect transfers across all test conditions.

\subsection{Intervention Gradient}

\begin{table}[h]
\centering
\begin{tabular}{lccc}
\toprule
\textbf{Displacement} & \textbf{Dead Baseline} & \textbf{PPO\_124 final} \\
\midrule
$\pm$0 px   & 0.0\% & 37.5\% \\
$\pm$8 px   & 0.0\% & 41.2\% \\
$\pm$15 px  & 0.0\% & \textbf{60.0\%} \\
$\pm$30 px  & 0.0\% & 50.0\% \\
$\pm$45 px  & 0.0\% & 31.2\% \\
$\pm$60 px  & 0.0\% & 25.0\% \\
\midrule
\textbf{AUC} & 0.000 & \textbf{0.421} \\
\bottomrule
\end{tabular}
\caption{Intervention gradient: reversal rate by displacement magnitude.
Dead baseline is a center-hold script that ignores the teleport.}
\label{tab:intervention}
\end{table}

Table~\ref{tab:intervention} shows the dose-response curve. The reversal rate
peaks at 60.0\% at $\pm$15px displacement---a moderate perturbation where the
ball is clearly in a new position but still within a physically plausible
range. The rate declines smoothly at larger displacements ($\pm$45--60px),
where the ball appears in physically impossible positions. This is the
signature of a reactive policy: the model recognizes the ball moved and
reverses direction to follow it, but when the displacement is too extreme, it
does not chase a ball that appears to violate physics. The dead baseline shows
0.0\% reversal at all magnitudes, confirming that the intervention test
correctly identifies non-reactive behavior.

The AUC of 0.421 is classified as STRONG and represents a clean dose-response
curve that memorized models never produce. Prior cursor models (PPO\_107--117)
showed 33--50\% reversal rates on the intervention probe but were confirmed
memorized by the split-watcher---those models had reactive \emph{distributions}
but memorized \emph{argmax modes}.

\subsection{Memorization Check}

Table~\ref{tab:memcheck} shows the memorization check results (deterministic
inference on clean Breakout, no proximity reward). PPO\_124 is the first model
in the project's history to sustain MULTIPLE\_SCRIPTS on det=True without
sticky action masking, with 10 of the last 12 checkpoints (14M--25M steps)
producing multiple unique scores. Earlier models achieving MULTIPLE\_SCRIPTS
did so only with sticky actions enabled (a known confound: sticky noise
produces 8--14 unique scores from a dead policy).

\begin{table}[h]
\centering
\begin{tabular}{lccc}
\toprule
\textbf{Checkpoint} & \textbf{Unique Scores} & \textbf{Best Score} &
\textbf{Verdict} \\
\midrule
1M  & 1 & 16  & SINGLE\_SCRIPT \\
5M  & 1 & 37  & SINGLE\_SCRIPT \\
10M & 1 & 78  & SINGLE\_SCRIPT \\
14M & 4 & 87  & MULTIPLE\_SCRIPTS \\
19M & 2 & 107 & SINGLE\_SCRIPT \\
25M & 4 & 93  & MULTIPLE\_SCRIPTS \\
\bottomrule
\end{tabular}
\caption{Memorization check: det=True unique scores on clean Breakout.}
\label{tab:memcheck}
\end{table}

\subsection{Continuation Training}

Extending training from 25M to 50M steps (PPO\_126, identical configuration)
did \emph{not} improve reactivity. The best checkpoint at 47.4M showed
layout-specific partial decoupling (RIGHT\_HALF: px\_corr=0.33, ALT=223),
but the final checkpoint at 50M converged to a single memorized script
(px\_corr=0.95 on all layouts, ALT=401 every game). Intervention AUC at 50M
was 0.327 (vs. 0.421 for PPO\_124 final). PPO's optimizer eventually finds a
script that maximizes the combined game + proximity objective; the proximity
reward delays but does not permanently prevent script convergence. This
suggests checkpoint selection---rather than training to convergence---is
critical for proximity-reward models.

%-------------------------------------------------------------------------------
\section{Discussion}

\subsection{Why Proximity Reward Works}

Every prior anti-memorization method in this project tried to change the
\emph{viability of scripts}: make the environment stochastic so scripts fail
(sticky actions, random ball perturbations), add penalties so repetitive
sequences are punished (entropy bonuses, trajectory entropy), or make the
environment geometrically complex (moving bumpers, brick randomization). PPO's
optimizer always found a script that survived the changes: timing-robust,
layout-conditioned, or noise-tolerant. The optimum was always a script---only
the shape changed.

Proximity reward changes \emph{what the optimum is}. A center-hold script---the
default PPO convergence point for Breakout---receives incidental proximity
bonus when the ball happens to pass near center. But a reactive tracking policy
receives the maximum bonus on every single descent frame, earning 3--5$\times$
more proximity reward over a full game. There is no clever script that can fake
being close to the ball. The optimization pressure is unambiguous: track the
ball, get more reward.

The bonus is dense (every frame) while the game's natural reward---brick
breaks---is sparse (every few seconds). Dense rewards provide better gradients
for credit assignment. The model does not need to discover the multi-step
causal chain ``track ball $\rightarrow$ hit ball $\rightarrow$ ball breaks
brick $\rightarrow$ score.'' It is told directly: being near the ball is good.
And because the scale is small (0.05 vs. 1.0--7.0 per brick), it does not
overwhelm the game reward. The model still cares about breaking bricks; it
learns that the way to break bricks is to be under the ball when it comes down.

\subsection{The Distribution-vs-Argmax Confound}

A critical methodological finding of this work is that \emph{every} commonly
used RL diagnostic---memorization check (unique score count), intervention
probe (distribution shift after state perturbation), SCAD probe (mutual
information between action and state), brick layout test (score
retention)---measures properties of the \emph{policy distribution}, not the
\emph{argmax action}. But evaluation uses the argmax. A model can maintain a
reactive-looking distribution (90\% LEFT when ball is left, 90\% RIGHT when
ball is right) while the argmax is always the same action regardless.

This was confirmed empirically: cursor-wrapper models (PPO\_107--117) showed
33--50\% reversal rates on the intervention probe and MULTIPLE\_SCRIPTS on
memcheck, but were definitively memorized under split-watcher verification
(perfect transfer on altered layouts). The intervention probe correctly
measures that the distribution $P(a \mid s)$ changes when the ball position
changes. It does not measure whether $\arg\max_a P(a \mid s)$ changes. For
models with high-confidence distributions (e.g., 90\% probability on the
top action), substantial probability mass can redistribute among the remaining
actions without ever crossing the argmax threshold.

The split-watcher is, to our knowledge, the only existing diagnostic that
directly measures argmax behavior by comparing paddle trajectories on different
initial states. We recommend it as a standard verification step for any claim
of reactive behavior in deterministic RL environments.

\subsection{Limitations}

\begin{enumerate}
    \item \textbf{Single game.} This work demonstrates the proximity reward
    approach only on Atari Breakout. The principle---directly reward the
    behavior you want---should generalize to other environments, but the
    specific implementation (reading RAM addresses for ball/paddle position)
    requires per-game engineering.
    \item \textbf{Scale sensitivity.} The scale of 0.05 was found to work; the
    sensitivity of reactivity to different scales (0.01, 0.10, 0.25) has not
    been systematically tested.
    \item \textbf{Finite training window.} PPO\_126's regression at 50M steps
    shows that the proximity reward delays but does not permanently prevent
    script convergence. The duration of the ``reactive window'' and whether it
    can be extended through scale scheduling remains unexplored.
    \item \textbf{No statistical significance testing.} Results are reported as
    raw counts and percentages without confidence intervals. The 0/240 perfect
    transfers result is qualitatively clear, but formal statistical treatment
    would strengthen the conclusions.
    \item \textbf{No ALE sticky evaluation protocol.} The standard ALE
    evaluation protocol (sticky actions enabled during evaluation) was not
    used here, as sticky actions are a known memorization masker. Clean
    evaluation (no sticky actions) provides the most direct test of argmax
    behavior.
\end{enumerate}

\subsection{Future Work}

The proximity reward approach opens several directions:

\begin{itemize}
    \item \textbf{Generalization to other Atari games.} The mechanism
    (dense reward for spatial alignment with a game entity) is Breakout-specific
    but the principle (reward the desired behavior directly) should transfer.
    Space Invaders: reward horizontal alignment with enemies. BeamRider: reward
    being out of the line of fire. Pong: reward paddle-ball proximity (trivially,
    since the game already does).
    \item \textbf{Scale scheduling.} Since PPO\_126 showed regression at 50M,
    could a scheduled decrease in proximity scale (start high to establish
    tracking, reduce to let game reward dominate) produce a stable reactive
    policy?
    \item \textbf{Combination with auxiliary supervision.} PPO\_102/103 showed
    that NatureCNN can track the ball at 1.9px MAE when supervised. Combining
    proximity reward (shapes the policy) with auxiliary ball-position
    supervision (shapes the features) might produce even stronger reactivity.
    \item \textbf{Formal theory.} Under what conditions does a dense shaped
    reward shift the global optimum of PPO's objective from a memorized script
    to a reactive policy? A theoretical treatment of the reward-shaping
    landscape for deterministic MDPs would be valuable.
\end{itemize}

%-------------------------------------------------------------------------------
\section{Conclusion}

After 124 controlled PPO experiments on Atari Breakout spanning every major
proposed anti-memorization method, we found that no environment-modification
approach produced a genuinely reactive argmax. Every diagnostic in the standard
RL toolbox measures the policy distribution, not the argmax---a confound that
led to false positive reactivity claims across multiple experimental
generations.

The fix was three lines of reward shaping: a dense proximity bonus that rewards
the paddle for horizontal closeness to the ball during descent. This directly
makes ball-tracking the reward-maximizing behavior, rather than trying to make
scripts non-viable through environment engineering. The resulting model
(PPO\_124) achieved 0/240 perfect transfers on the split-watcher, 100\%
no-timing ALT retention, and an intervention gradient AUC of 0.421---the first
verified reactive PPO argmax on Atari Breakout.

The core lesson is simple: PPO optimizes whatever you put in the reward
function. In a deterministic environment, that optimum is a script---unless you
change the reward to make it something else. Reward what you want, don't
penalize what you don't want.

%-------------------------------------------------------------------------------
\bibliographystyle{unsrt}
\begin{thebibliography}{99}

\bibitem{schulman2017ppo}
J.~Schulman, F.~Wolski, P.~Dhariwal, A.~Radford, and O.~Klimov.
\newblock Proximal policy optimization algorithms.
\newblock {\em arXiv preprint arXiv:1707.06347}, 2017.

\bibitem{bellemare2013ale}
M.~G.~Bellemare, Y.~Naddaf, J.~Veness, and M.~Bowling.
\newblock The arcade learning environment: An evaluation platform for general
  agents.
\newblock {\em Journal of Artificial Intelligence Research}, 47:253--279, 2013.

\bibitem{machado2018revisiting}
M.~C.~Machado, M.~G.~Bellemare, E.~Talvitie, J.~Veness, M.~Hausknecht, and
  M.~Bowling.
\newblock Revisiting the arcade learning environment: Evaluation protocols and
  open problems for general agents.
\newblock {\em Journal of Artificial Intelligence Research}, 61:523--562, 2018.

\bibitem{zhang2018overfitting}
C.~Zhang, O.~Vinyals, R.~Munos, and S.~Bengio.
\newblock A study on overfitting in deep reinforcement learning.
\newblock {\em arXiv preprint arXiv:1804.06893}, 2018.

\bibitem{justesen2018illuminating}
N.~Justesen, R.~R.~Torrado, P.~Bontrager, A.~Khalifa, J.~Togelius, and
  S.~Risi.
\newblock Illuminating generalization in deep reinforcement learning through
  procedural level generation.
\newblock In {\em NeurIPS Workshop on Deep Reinforcement Learning}, 2018.

\bibitem{cobbe2020procgen}
K.~Cobbe, C.~Hesse, J.~Hilton, and J.~Schulman.
\newblock Leveraging procedural generation to benchmark reinforcement learning.
\newblock In {\em International Conference on Machine Learning}, 2020.

\bibitem{ecoffet2021goexplore}
A.~Ecoffet, J.~Huizinga, J.~Lehman, K.~O.~Stanley, and J.~Clune.
\newblock First return, then explore.
\newblock {\em Nature}, 590(7847):580--586, 2021.

\bibitem{ng1999shaping}
A.~Y.~Ng, D.~Harada, and S.~Russell.
\newblock Policy invariance under reward transformations: Theory and application
  to reward shaping.
\newblock In {\em International Conference on Machine Learning}, 1999.

\bibitem{jaderberg2017auxiliary}
M.~Jaderberg, V.~Mnih, W.~M.~Czarnecki, T.~Schaul, J.~Z.~Leibo, D.~Silver,
  and K.~Kavukcuoglu.
\newblock Reinforcement learning with unsupervised auxiliary tasks.
\newblock In {\em International Conference on Learning Representations}, 2017.

\bibitem{raffin2021sb3}
A.~Raffin, A.~Hill, A.~Gleave, A.~Kanervisto, M.~Ernestus, and N.~Dormann.
\newblock Stable-baselines3: Reliable reinforcement learning implementations.
\newblock {\em Journal of Machine Learning Research}, 22(268):1--8, 2021.

\bibitem{mnih2015human}
V.~Mnih, K.~Kavukcuoglu, D.~Silver, A.~A.~Rusu, J.~Veness, M.~G.~Bellemare,
  A.~Graves, M.~Riedmiller, A.~K.~Fidjeland, G.~Ostrovski, et~al.
\newblock Human-level control through deep reinforcement learning.
\newblock {\em Nature}, 518(7540):529--533, 2015.

\end{thebibliography}

\end{document}
